%% ma: L10-B01-0030
%% lop: 10
%% chuong: 1
%% bai: 1
%% dang: 10.1.2
%% loai: DS
%% muc: [NB, TH, NB, NB]
%% dap_an: "ĐSĐS"
%% nguon: "Ảnh giáo viên gửi 10/10/2026 (ThS. Nguyễn Hoàng Việt, Chương 1 Mệnh đề), B-Đề 2, Phần 2 Câu 13"
%% kien_thuc: "Bài 1 (mệnh đề, mệnh đề chứa biến)"
%% trang_thai: cho-duyet
%% ly_do: "Dạng toán mới đề xuất 10.1.2: xét tính đúng sai của mệnh đề (số học, đại số, hình học)"
\begin{ex}
Xét tính đúng, sai của mỗi mệnh đề sau.
\choiceTF
{\True $15$ không là số nguyên tố.}
{Một tứ giác là hình thoi khi và chỉ khi nó có hai đường chéo vuông góc với nhau.}
{\True $5+19=24$.}
{$6+81=25$.}
\loigiai{a) $15=3\cdot5$. Đúng. b) Sai: tứ giác có hai đường chéo vuông góc chưa chắc là hình thoi (ví dụ hình diều). c) Đúng. d) $6+81=87\ne25$. Sai.}
\end{ex}
